\documentclass[11pt,a4paper]{article}
\usepackage[margin=2.4cm]{geometry}
\usepackage{amsmath,amssymb,amsthm,mathtools}
\usepackage{booktabs}
\usepackage{enumitem}
\usepackage{xcolor}
\usepackage{hyperref}
\usepackage{fontspec}
\usepackage{polyglossia}
\setmainlanguage{russian}
\setotherlanguage{english}
\setmainfont{Times New Roman}
\setmonofont{Consolas}[Scale=0.90]

\hypersetup{
  colorlinks=true,
  linkcolor=blue!70!black,
  citecolor=blue!70!black,
  urlcolor=blue!70!black,
  pdftitle={Основания дискретной квантовой механики (ДКМ): Единственность размерности 1 + 2^d = d^2 (d = 3), симплектическая динамика Верле на Z_512, правило Борна и постоянная тонкой структуры в Lean 4},
  pdfauthor={Каратаев Дмитрий Александрович}
}

\theoremstyle{plain}
\newtheorem{theorem}{Теорема}[section]
\newtheorem{lemma}[theorem]{Лемма}
\newtheorem{proposition}[theorem]{Предложение}
\newtheorem{corollary}[theorem]{Следствие}
\theoremstyle{definition}
\newtheorem{definition}[theorem]{Определение}
\newtheorem{remark}[theorem]{Замечание}

\title{\textbf{Основания дискретной квантовой механики (ДКМ):\\ Единственность размерности $1 + 2^d = d^2$ ($d = 3$), симплектическая динамика Верле на $\mathbb{Z}_{512}$, правило Борна и постоянная тонкой структуры $\alpha_{\mathrm{fs}}^{-1}$ в Lean~4}}
\author{Д.\,А.~Каратаев\thanks{AI-Guiders Research (\url{https://github.com/AI-Guiders}). Формальная верификация в Lean~4 и аналитическая разработка выполнены при участии ИИ-партнёра \emph{Луч} (\emph{Luch}, AI-Guiders). Работа распространяется на условиях лицензии строгого взаимного Copyleft: текст и таблицы --- Creative Commons Attribution-ShareAlike 4.0 International (\texttt{CC-BY-SA-4.0}); код формальных доказательств Lean~4 --- GNU General Public License v3.0 or later (\texttt{GPL-3.0-or-later}).}}
\date{2026}

\begin{document}
\maketitle

\begin{abstract}
В настоящей работе по основаниям дискретной квантовой механики (ДКМ) строится обратимый локальный клеточный автомат над двусортной алгебраической системой $\langle\mathcal{M},\mathcal{T};\,\pi_{M\to T},\,\iota_{T\to M}\rangle$ в смысле теории моделей А.\,И.~Мальцева (\emph{«Алгебраические системы»}), где \textbf{дискретный $M$-носитель} есть бесконечное счётное индуктивное семейство целочисленных фазовых колец $\mathbb{Z}_{4Q}$ на решётке $\mathbb{Z}^d$, а \textbf{непрерывный $T$-носитель} --- континуальный сорт наблюдаемых полей, траекторий и интегралов Римана. Физическая неделимость элементарного объёма (частица занимает весь элементарный объём ячейки целиком и не может занимать его частично, что даёт $1$ неделимое состояние собственной занятости центра) и плотная упаковка всех ячеек без пустот вдоль числа осей $d$ ($2^d$ сомкнутых знаковых октантов) порождают $1 + 2^d$ неделимых состояний ветвления ячейки. Их взаимно-однозначная замкнутость с $d^2$ упорядоченными парами осевых переходов даёт диофантово уравнение Каталана--Михайлеску $1 + 2^d = d^2$, однозначно выделяющее размерность пространства $d = 3$, разрядность регистра $9$ бит, модуль фазового кольца $N_{\mathrm{ring}}(3) = 512$, вакуумный квант амплитуды $\mu_{\mathrm{vac}}(3) = 1$, четверть-период $Q_\square(3) = 128$, потолок Паули $\Phi_{\mathrm{Pauli}}(3) = 256$ и шаг биссектрисы $K_P(3) = 64$.

Доказано, что единичный симплектический якобиан, обратимость во времени и глобальная калибровочная инвариантность однозначно фиксируют двухшаговый закон эволюции Верле на $72$-битном регистре ячейки $B_2(3) = 72$, а изотропия второго момента и замкнутость планарных $4$-циклов в кубе Мура однозначно выделяют $12$-вершинный кубооктаэдр $O_2(3)$ ($13$-узловое причинное ядро $B_1(3)$). Деление с остатком разности фаз по модулю $512$ расщепляет динамику на волновой остаток и корпускулярный перенос, детерминированно выводя правило Борна $P(\mu) = \mu^2/512$, дискретные уравнения Максвелла--Фарадея--Гаусса, порог Гейзенберга $\Delta\varphi_{\mathrm{disc}}(3) = 41$ (подпороговый коридор без тепла Ландауэра и закон проекционной слепоты $64\times 81 = 72^2 = 5\,184$), явные структурные константы коммутаторов калибровочной алгебры Ли Стандартной модели $\mathfrak{u}(1)\oplus\mathfrak{su}(2)\oplus\mathfrak{su}(3)$ размерности $1 + 3 + 8 = 12$ на кликовом комплексе $\mathcal{C}_{13}$ и четырёхпетлевой вывод постоянной тонкой структуры в поле $\mathbb{Q}(\sqrt{2})$ без трансцендентного числа~$\pi$: $\alpha_{\mathrm{fs}}^{-1}(3) = 137{,}035\,999\,187\dots$ (отклонение $0{,}47\,\text{ст.\ откл.}$ от CODATA-2022). Все теоремы формально верифицированы в \textbf{Lean~4} (\texttt{0 sorryAx}, \texttt{0 Classical.choice}).
\end{abstract}

\section{Постановка задачи: почему физическая среда без внешней памяти обязана быть дискретной, трёхмерной и описываться кубооктаэдром \texorpdfstring{$O_2(3)$}{O2(3)} и динамикой Верле}
\label{sec:axioms-uniqueness}

В классической теоретической физике пространство-время традиционно вводится как готовый гладкий континуум $\mathbb{R}^{d,1}$, поверх которого задаются дифференциальные уравнения поля и эмпирические константы связи. Такой подход неявно опирается на идеализацию \emph{внешнего наблюдателя}: предполагается, что состояние сколь угодно малого элемента объёма может быть записано вещественными числами бесконечной разрядности на некотором внешнем носителе за пределами самой физической системы.

Отказ от этой идеализации приводит к фундаментальному физическому вопросу: \emph{каким образом должна быть устроена среда, которая вычисляет собственную эволюцию шаг за шагом, не обращаясь ни к какой внешней оперативной памяти за своими пределами?}

Поскольку любой измерительный прибор и любой наблюдатель сами являются физическими подсистемами внутри той же самой среды, отсутствие внешней памяти влечёт прямую четырёхшаговую цепочку физических следствий:
\begin{enumerate}[label=\textbf{Шаг \arabic*.},leftmargin=*]
\item \textbf{Дискретность на уровне элементарного объёма и двусортная структура $\langle\mathcal{M},\mathcal{T};\,\pi_{M\to T},\,\iota_{T\to M}\rangle$ в смысле А.\,И.~Мальцева~\cite{Malcev1970}.}
  Частица (возбуждение) занимает весь элементарный объём ячейки целиком и не может занимать его частично. В среде без внешней памяти состояние одного оборота каждого узла $\mathbf{x}\in\mathbb{Z}^d$ на такте $t\in\mathbb{Z}$ кодируется конечным двоичным словом фиксированной разрядности $\mathrm{mod\_bits}(d)$ на циклическом кольце вычетов $\mathbb{Z}_{N_{\mathrm{ring}}(d)}$, где
  \begin{equation}
    \label{eq:nring-power-of-two}
    N_{\mathrm{ring}}(d) \;=\; 2^{\mathrm{mod\_bits}(d)}.
  \end{equation}
  В соответствии с теорией алгебраических систем А.\,И.~Мальцева~\cite{Malcev1970} система задаётся двумя сортами (\textbf{носителями}):
  \begin{itemize}[leftmargin=*]
  \item \textbf{дискретный $M$-носитель} $\mathcal{M}$ --- бесконечное счётное индуктивное семейство целочисленных фазовых колец $\mathbb{Z}_{4Q}$ по всем масштабам четверть-периода~$Q$, всем числам витков $q\in\mathbb{Z}$ и всем узлам решётки $\mathbb{Z}^d$ (с однооборотным фактор-регистром ячейки $\mathbb{Z}_{512}$ на масштабе $K=1$, $Q_\square(3)=128$);
  \item \textbf{непрерывный $T$-носитель} $\mathcal{T}$ --- континуальный сорт наблюдаемых полей, траекторий и интегралов Римана, связанный с~$\mathcal{M}$ проективным подъёмом Вейерштрасса $\pi_{M\to T}$ и линейной ретракцией $\iota_{T\to M}$.
  \end{itemize}

\item \textbf{Единственность трёхмерности пространства ($d = 3$), кольцо $\mathbb{Z}_{512}$ и квант $\mu_{\mathrm{vac}} = 1$.}
  Рассмотрим элементарную ячейку в узле $\mathbf{x}\in\mathbb{Z}^d$:
  \begin{itemize}[leftmargin=*]
  \item \emph{Неделимость элементарного объёма ($+1$ состояние центра):} поскольку частица занимает весь элементарный объём ячейки целиком, у центра ячейки существует ровно одно неделимое состояние собственной занятости $\mathbf{0} = (0,\dots,0)$ (\texttt{centralOccupancyState d} в \texttt{lean/MtCarrier/Spec.lean});
  \item \emph{Плотное заполнение пространства без пустот ($2^d$ сомкнутых октантов):} вдоль каждой из $d$ ортогональных осей существуют ровно два сомкнутых противоположных направления $\{-1,+1\}$, прямое произведение которых по $d$ осям образует $|\texttt{signedOctants}(d)| = 2^d$ знаковых октантов $\{-1,+1\}^d$. Вместе с состоянием центра это даёт
    \begin{equation}
      \label{eq:cell-branching-states}
      |\texttt{cellBranchingStates}(d)| \;=\; 1 + 2^d
    \end{equation}
    неделимых состояний ветвления ячейки;
  \item \emph{Матрица осевых переходов внутри ячейки ($d^2$ пар):} передача состояния между $d$ ортогональными осями внутри ячейки задаётся упорядоченными парами осей $(i,j)\in\{0,\dots,d-1\}^2$, число которых равно
    \begin{equation}
      \label{eq:axis-transition-pairs}
      |\texttt{axisTransitionPairs}(d)| \;=\; \mathrm{mod\_bits}(d) \;=\; d^2.
    \end{equation}
  \end{itemize}
  Биективная замкнутость между состояниями ветвления ячейки~\eqref{eq:cell-branching-states} и парами её осевых переходов~\eqref{eq:axis-transition-pairs} (\texttt{catalan\_lock\_eq\_branching\_axis\_cardinalities}) даёт диофантово уравнение Каталана--Михайлеску:
  \begin{equation}
    \label{eq:catalan-lock}
    1 + 2^d \;=\; d^2.
  \end{equation}

\item \textbf{Выделение $12$-вершинного кубооктаэдра $O_2(3)$ среди орбит куба Мура.}
  В размерности $d = 3$ куб Мура $\{-1,0,+1\}^3$ вокруг $(0,0,0)\in\mathbb{Z}^3$ содержит $3^3 - 1 = 26$ ненулевых соседей, распадающихся по квадрату длины связи $|e|_2^2 = e_x^2 + e_y^2 + e_z^2$ на три $O_h$-орбиты: $O_1(3)$ ($|e|_2^2 = 1$, $6$ вершин, где смешанный четвёртый момент равен нулю и оси не сцеплены), $O_3(3)$ ($|e|_2^2 = 3$, $8$ угловых вершин без планарных $4$-циклов в координатных плоскостях) и $O_2(3)$ ($|e|_2^2 = 2$, $12$ рёберных диагональных соседей --- кубооктаэдр).

\item \textbf{Выделение двухшаговой симплектической схемы Верле с коэффициентами $(2,-1)$.}
  Обратимая во времени двухшаговая рекурсия на кольце $\mathbb{Z}_{512}$ имеет вид
  \begin{equation}
    \label{eq:general-two-step}
    \varphi(t+1) \;\equiv\; \alpha\,\varphi(t) + \beta\,\varphi(t-1) + \Phi_{\mathrm{gate}} \pmod{512}.
  \end{equation}
  Сохранение симплектической площади $\det J = -\beta = +1$ фиксирует $\beta = -1$, а инвариантность при глобальном сдвиге фазы $\varphi \mapsto \varphi + c$ требует $\alpha + \beta = 1$, откуда $\alpha = 2$.
\end{enumerate}

\begin{theorem}[Единственность размерности $1 + 2^d = d^2$ ($d=3$) и константы кольца $\mathbb{Z}_{512}$]
\label{thm:d3-uniqueness}
Для $d\in\mathbb{N}$ диофантово уравнение самозамкнутости регистра~\eqref{eq:catalan-lock} выполняется тогда и только тогда, когда $d = 3$ (\texttt{catalan\_lock\_iff\_d3}, \texttt{canonical\_d3\_constants}), однозначно фиксируя без свободных параметров:
\begin{align}
  d &= 3, \qquad N_{\mathrm{sq}}(3) = 2d = 6, \qquad N_{\mathrm{tri}}(3) = 2^d = 8, \label{eq:d3-consts-1} \\
  \mathrm{mod\_bits}(3) &= d^2 = 9, \qquad N_{\mathrm{ring}}(3) = 2^9 = 512, \qquad \mu_{\mathrm{vac}}(3) = d^2 - 2^d = 1, \label{eq:d3-consts-2} \\
  Q_\square(3) &= 128, \qquad \Phi_{\mathrm{Pauli}}(3) = 256, \qquad K_P(3) = 64. \label{eq:d3-consts-3}
\end{align}
\end{theorem}

\begin{proof}
Прямая проверка уравнения $1 + 2^d = d^2$ при $d \in \{0, 1, 2, 3, 4\}$ даёт:
\begin{align*}
  d = 0 &: \quad 1 + 2^0 = 2 \ne 0^2, \\
  d = 1 &: \quad 1 + 2^1 = 3 \ne 1^2, \\
  d = 2 &: \quad 1 + 2^2 = 5 \ne 2^2, \\
  d = 3 &: \quad 1 + 2^3 = 9 = 3^2, \\
  d = 4 &: \quad 1 + 2^4 = 17 \ne 4^2.
\end{align*}
Для всех $d \ge 4$ докажем по индукции неравенство $2^d \ge d^2$. При $d = 4$ имеем $2^4 = 16 = 4^2$. Пусть $2^d \ge d^2$ при некотором $d \ge 4$; тогда
\begin{equation*}
  2^{d+1} \;=\; 2\cdot 2^d \;\ge\; 2d^2 \;=\; d^2 + d\cdot d \;\ge\; d^2 + 4d \;>\; d^2 + 2d + 1 \;=\; (d + 1)^2,
\end{equation*}
откуда $1 + 2^d > d^2$ для всех $d \ge 4$. Следовательно, $d = 3$ --- единственный корень в $\mathbb{N}$. Подстановка $d = 3$ даёт константы~\eqref{eq:d3-consts-1}--\eqref{eq:d3-consts-3}.
\end{proof}

\begin{theorem}[Единственность симплектического закона Верле и оболочки кубооктаэдра $O_2(3)$]
\label{thm:d3-verlet-cuboct}
Справедливы следующие теоремы единственности (\texttt{verlet\_coefficients\_uniqueness\_from\_a13}, \texttt{polar\_step\_reverse\_exact}, \texttt{cuboctahedron\_unique\_isotropic\_shell\_d3}):
\begin{enumerate}[label=\textbf{(\arabic*)},leftmargin=*]
\item Среди двухшаговых линейных рекурсий~\eqref{eq:general-two-step} условия единичного якобиана $\det J = -\beta = 1$ и инвариантности относительно сдвига фазы $\alpha + \beta = 1$ однозначно фиксируют $(\alpha, \beta) = (2, -1)$:
  \begin{equation}
    \label{eq:verlet-law}
    \Delta^2\varphi(t) \;:=\; \varphi(t+1) - 2\varphi(t) + \varphi(t-1) \;\equiv\; \Phi_{\mathrm{gate}} \pmod{512},
  \end{equation}
  задавая обратимую инволюцию на $72$-битном двухтактовом фазовом пространстве ячейки $B_2(3) = 2\times 36 = 72\text{ бит}$.
\item В трёхмерном кубе Мура $\{-1,0,+1\}^3$ единственной минимальной оболочкой, обеспечивающей изотропию второго момента
  \begin{equation}
    \label{eq:o2-second-moment}
    \sum_{e\in O_2(3)} e_\alpha e_\beta \;=\; 8\,\delta_{\alpha\beta},
  \end{equation}
  замкнутость $4$-циклов в координатных плоскостях $(xy,yz,zx)$ и точное воспроизведение оператора Лапласа на квадратичных многочленах, является $12$-вершинный кубооктаэдр
  \begin{equation}
    \label{eq:o2-def}
    O_2(3) \;=\; \bigl\{\,e\in\{-1,0,+1\}^3 \;\big|\; e_x^2 + e_y^2 + e_z^2 = 2\,\bigr\}
  \end{equation}
  с $13$-узловым причинным ядром $B_1(3) = \{\mathbf{0}\}\cup O_2(3)$ и законом роста координационных слоёв $S_k(3) = 10k^2 + 2$ ($12, 42, 92, 162$).
\end{enumerate}
\end{theorem}

\begin{proof}
В пункте~(1) матрица перехода одного шага $(\varphi(t), \varphi(t-1)) \mapsto (\alpha\varphi(t) + \beta\varphi(t-1), \varphi(t))$ имеет определитель $\det J = \alpha\cdot 0 - \beta\cdot 1 = -\beta$. Условие $-\beta = 1$ даёт $\beta = -1$, а условие $\alpha + \beta = 1$ даёт $\alpha = 1 - (-1) = 2$. Перестановка $(\varphi(t+1), \varphi(t)) \leftrightarrow (\varphi(t-1), \varphi(t))$ сохраняет сумму $\varphi(t+1) + \varphi(t-1) \equiv 2\varphi(t) + \Phi_{\mathrm{gate}} \pmod{512}$, что доказывает обратимость во времени.

В пункте~(2) куб Мура $\{-1,0,+1\}^3 \setminus \{\mathbf{0}\}$ разбивается на три орбиты $O_1(3)$ ($6$ вершин), $O_2(3)$ ($12$ вершин) и $O_3(3)$ ($8$ вершин). На $O_1(3)$ смешанный момент $\sum e_x^2 e_y^2 = 0$, поэтому оси не сцеплены; на $O_3(3)$ у всех $8$ вершин все три координаты ненулевые ($|e_i| = 1$), поэтому в координатных плоскостях $e_z = 0$, $e_x = 0$, $e_y = 0$ нет ни одной вершины. Только $O_2(3)$ содержит три экваториальных квадрата по $4$ вершины в координатных плоскостях и имеет моменты $\sum e_i^2 = 8$, $\sum_{i\ne j} e_i^2 e_j^2 = 4$, обеспечивая~\eqref{eq:o2-second-moment} и тождество $\frac{1}{4}\sum_{e\in O_2(3)}(f(e) - f(\mathbf{0})) = \Delta f(\mathbf{0})$ на всех квадратичных формах.
\end{proof}

\section{Правило Борна, порог Гейзенберга \texorpdfstring{$\Delta\varphi_{\mathrm{disc}}=41$}{dphi=41}, уравнения Максвелла и закон проекционной слепоты \texorpdfstring{$72^2=5\,184$}{72^2=5184}}
\label{sec:born-maxwell-heisenberg}

\begin{theorem}[Корпускулярно-волновое расщепление, правило Борна $P(\mu)=\mu^2/512$ и дискретные уравнения Максвелла--Гаусса--Фарадея]
\label{thm:born-maxwell}
При $d=3$ на кольце $\mathbb{Z}_{512}$ и оболочке $O_2(3)$ выполняются следующие законы (\texttt{quotient\_remainder\_512\_spec}, \texttt{bornClickCountOp}, \texttt{d1\_d0\_zero}, \texttt{discrete\_gauss\_cube\_divergence\_zero}):
\begin{enumerate}[label=\textbf{(\arabic*)},leftmargin=*]
\item Всякая разность фаз $\delta = \varphi_y - \varphi_x \in (-512, 512)$ однозначно расщепляется по формуле
  \begin{equation}
    \label{eq:phase-qr-split}
    \delta \;=\; 512\,q(\delta) + r(\delta)
  \end{equation}
  на волновой остаток $r(\delta) \in (-256, 256]$ и корпускулярный перенос $q(\delta) \in \{-1, 0, +1\}$. Для волны амплитуды $\mu \in \{0, \dots, 16\}$ ($\Phi_{\mathrm{wave}} = \mu^2 \le 256$) на детекторе с равномерной фоновой фазой $\varphi_{\mathrm{bg}} \in \{0, \dots, 511\}$ число фоновых состояний, вызывающих срабатывание переноса $q(\varphi_{\mathrm{bg}} + \mu^2) = 1$, тождественно равно квадрату амплитуды:
  \begin{equation}
    \label{eq:born-rule-exact}
    N_{\mathrm{clicks}}(\mu) \;\equiv\; \mu^2 \;=\; |\psi|^2,
    \qquad
    P(\mu) \;=\; \frac{N_{\mathrm{clicks}}(\mu)}{512} \;=\; \frac{\mu^2}{512}.
  \end{equation}
\item Вокруг любой замкнутой $4$-цикловой грани $f$ циркуляция волновых остатков квантована числом витков $w_f = -\sum_{e\in\partial f} q_e \in \mathbb{Z}$, а полный поток вихрей через $6$ граней любого замкнутого куба $C_3$ тождественно равен нулю:
  \begin{align}
    \sum_{e\in\partial f} r_e &= -512\sum_{e\in\partial f} q_e \;=\; 512\,w_f, \label{eq:maxwell-faraday-circulation} \\
    \sum_{f\in\partial C_3} w_f &\equiv 0. \label{eq:maxwell-gauss-monopole-free}
  \end{align}
\end{enumerate}
\end{theorem}

\begin{proof}
В пункте~(1) при сложении входного фазового импульса $\Phi_{\mathrm{wave}} = \mu^2 \in [0, 256]$ с фоновой фазой детектора $\varphi_{\mathrm{bg}} \in \{0, 1, \dots, 511\}$ перенос через границу кольца $\varphi_{\mathrm{bg}} + \mu^2 \ge 512$ происходит в точности для $\mu^2$ значений $\varphi_{\mathrm{bg}} \in \{512 - \mu^2, \dots, 511\}$, откуда $N_{\mathrm{clicks}}(\mu) = \mu^2$ и $P(\mu) = \mu^2 / 512$.

В пункте~(2) сумма точных разностей $\delta_e = \varphi_y - \varphi_x$ вдоль замкнутого контура $\partial f$ телескопически сокращается в нуль: $\sum_{e\in\partial f} \delta_e = 0$. Подставляя~\eqref{eq:phase-qr-split}, получаем $\sum_{e\in\partial f} r_e + 512\sum_{e\in\partial f} q_e = 0$, что даёт~\eqref{eq:maxwell-faraday-circulation}. При суммировании циркуляций $w_f$ по $6$ ориентированным граням замкнутого куба $\partial C_3$ каждое из $12$ рёбер куба проходится дважды во взаимно противоположных направлениях, что тождественно обнуляет сумму~\eqref{eq:maxwell-gauss-monopole-free}.
\end{proof}

\begin{theorem}[Порог Гейзенберга $\Delta\varphi_{\mathrm{disc}}(3)=41$, подпороговый коридор без тепла Ландауэра и проекционная слепота $72^2=5\,184$]
\label{thm:born-heisenberg-maxwell}
Полный фазовый бюджет $13$-узлового ядра $B_1(3)$ со швом грани $\Delta\varphi_{\mathrm{seam}}(3) = \lfloor 256/12\rfloor = 21$ делится нацело на число узлов ядра $N_\varphi(3) = 13$ исключительно при $d = 3$ (в диапазоне $1 \le d \le 8$, \texttt{heisenberg\_divisibility\_lock\_unique\_d3\_up\_to\_8}, \texttt{a6\_a8\_gate\_pauli\_ceiling\_256\_and\_heisenberg\_deadband\_41}):
\begin{equation}
  \label{eq:heisenberg-533-factorization}
  N_{\mathrm{ring}}(3) + \Delta\varphi_{\mathrm{seam}}(3) \;=\; 512 + 21 \;=\; 533 \;=\; 13 \times 41 \;=\; N_\varphi(3)\,\Delta\varphi_{\mathrm{disc}}(3),
\end{equation}
фиксируя порог излучения $\Delta\varphi_{\mathrm{disc}}(3) = 41$ и число излучённых квантов
\begin{equation}
  \label{eq:heisenberg-quantum-count}
  \hat{n}_E(\Phi) \;=\; \bigl\lfloor |\Phi| / 41 \bigr\rfloor \;\in\; \{0, 1, 2, 3, 4, 5, 6\}
  \qquad (|\Phi| \le 256).
\end{equation}
В подпороговом коридоре $|\Phi| \le 40$, содержащем $|\mathcal{E}_{\mathrm{Heis}}(3)| = 2\cdot 41 - 1 = 81 = 3^4$ состояний, логические операции протекают обратимо без излучения ($\hat{n}_E = 0$), а макроскопический $T$-прибор склеивает
\begin{equation}
  \label{eq:projection-blindness-5184}
  K_P(3)\cdot |\mathcal{E}_{\mathrm{Heis}}(3)| \;=\; 64 \times 81 \;=\; 5\,184 \;=\; 72^2 \;=\; \bigl(B_2(3)\bigr)^2
\end{equation}
детерминированных микросостояний $M$-носителя в одно $T$-макросостояние.
\end{theorem}

\begin{proof}
При $d = 3$ имеем $N_\varphi(3) = 1 + |V_{O_2}(3)| = 1 + 12 = 13$ и $\Delta\varphi_{\mathrm{seam}}(3) = \lfloor 256 / 12\rfloor = 21$, откуда $512 + 21 = 533 = 13 \times 41$. Поскольку $|\Phi| \le \Phi_{\mathrm{Pauli}}(3) = 256$ и $6 \times 41 = 246 \le 256 < 287 = 7 \times 41$, целочисленное частное $\lfloor |\Phi|/41\rfloor$ принимает значения от $0$ до $6$, причём $\lfloor |\Phi|/41\rfloor = 0 \iff |\Phi| \le 40$ ($40$ отрицательных, $1$ нулевое и $40$ положительных состояний, всего $81$ состояние). Умножение $81$ подпорогового состояния на шаг октантной биссектрисы $K_P(3) = 64$ даёт $64 \times 81 = 8^2 \times 9^2 = 72^2 = 5\,184$, что совпадает с квадратом двухтактовой разрядности ячейки $B_2(3) = 2 \times 4 \times 9 = 72$.
\end{proof}

\section{Кликовый комплекс \texorpdfstring{$\mathcal{C}_{13}$}{C13}, коммутаторы алгебры Ли \texorpdfstring{$\mathfrak{u}(1)\oplus\mathfrak{su}(2)\oplus\mathfrak{su}(3)$}{u(1)+su(2)+su(3)} и четырёхпетлевой вывод постоянной тонкой структуры \texorpdfstring{$\alpha_{\mathrm{fs}}^{-1}$}{alpha\_fs\^{}(-1)}}
\label{sec:sm-alpha}

\begin{theorem}[Кликовый комплекс $\mathcal{C}_{13}$, матричные коммутаторы алгебры Ли $\mathfrak{u}(1)\oplus\mathfrak{su}(2)\oplus\mathfrak{su}(3)$, гауссово простое число Кулона $\widehat{M}(3)=97$ и вывод постоянной тонкой структуры $\alpha_{\mathrm{fs}}^{-1}$]
\label{thm:sm-alpha}
На $13$-узловом причинном ядре $B_1(3) = \{\mathbf{0}\}\cup O_2(3)$ выполняются следующие структурные законы (\texttt{gaussian\_coulomb\_lock\_iff\_d3}, \texttt{alpha\_fs\_exact\_z\_sqrt2\_canonical\_form\_and\_floor\_137}, \texttt{sm\_and\_cosmology\_closed\_rational\_operators}):
\begin{enumerate}[label=\textbf{(\arabic*)},leftmargin=*]
\item \textbf{Симплициальная топология комплекса $\mathcal{C}_{13}$ и явная реализация коммутаторов Ли $\mathfrak{u}(1)\oplus\mathfrak{su}(2)\oplus\mathfrak{su}(3)$:}
  числа непустых клик графа $B_1(3)$ по размерностям равны $(K_1, K_2, K_3, K_4, K_{\ge 5}) = (13, 36, 32, 8, 0)$ с эйлеровой характеристикой
  \begin{equation}
    \label{eq:c13-euler-char}
    \chi(\mathcal{C}_{13}) \;=\; 13 - 36 + 32 - 8 \;=\; 1 \;=\; \mu_{\mathrm{vac}}(3).
  \end{equation}
  На трёх координатных плоскостях $(yz, zx, xy)$ ядра $B_1(3)$ действуют три семейства целочисленных $3\times 3$ матричных генераторов с коммутатором Ли $[A, B] = AB - BA$:
  \begin{itemize}[leftmargin=*]
  \item $1$ центральный генератор глобального фазового сдвига $\mathfrak{u}(1)$: единичная матрица $Y_0 = I_3$, коммутирующая со всеми пространственными переходами ($[Y_0, L_k] = 0$, $[Y_0, E_{ab}] = 0$);
  \item $d = 3$ кососимметрических генератора поворотов на $90^\circ$ в трёх экваториальных квадратных $4$-циклах $(yz, zx, xy)$ оболочки $O_2(3)$:
    \begin{equation}
      \label{eq:su2-generators-matrices}
      L_1 = \begin{pmatrix} 0 & 0 & 0 \\ 0 & 0 & -1 \\ 0 & 1 & 0 \end{pmatrix}, \qquad
      L_2 = \begin{pmatrix} 0 & 0 & 1 \\ 0 & 0 & 0 \\ -1 & 0 & 0 \end{pmatrix}, \qquad
      L_3 = \begin{pmatrix} 0 & -1 & 0 \\ 1 & 0 & 0 \\ 0 & 0 & 0 \end{pmatrix},
    \end{equation}
    образующих замкнутую алгебру Ли $\mathfrak{so}(3) \cong \mathfrak{su}(2)$ со структурными константами:
    \begin{equation}
      \label{eq:su2-commutators}
      [L_1, L_2] \;=\; L_3, \qquad
      [L_2, L_3] \;=\; L_1, \qquad
      [L_3, L_1] \;=\; L_2;
    \end{equation}
  \item $K_4(3) = d^2 - 1 = 8$ бесследовых генераторов Картана--Вейля алгебры $\mathfrak{sl}_3 \cong \mathfrak{su}(3)$ на триаде трёх координатных плоскостей вершин $8$ тетраэдрических $4$-клик $K_4$ ($6$ внедиагональных матриц межплоскостных переходов $E_{ab}$ при $a \ne b \in \{1,2,3\}$ и $2$ бесследовых диагональных генератора Картана $H_1 = E_{11} - E_{22}$, $H_2 = E_{22} - E_{33}$ с $\mathrm{Tr}(H_1) = \mathrm{Tr}(H_2) = 0$), удовлетворяющих структурным соотношениям Ли:
    \begin{align}
      [E_{ab}, E_{bc}] &= E_{ac} \qquad (a \ne c), \label{eq:su3-comm-step} \\
      [E_{12}, E_{21}] &= H_1, \qquad [E_{23}, E_{32}] \;=\; H_2, \qquad [H_1, H_2] \;=\; 0, \label{eq:su3-comm-cartan}
    \end{align}
  \end{itemize}
  что даёт полную калибровочную размерность Стандартной модели:
  \begin{equation}
    \label{eq:sm-gauge-dim-12}
    \dim\bigl(\mathfrak{u}(1)\oplus\mathfrak{su}(2)\oplus\mathfrak{su}(3)\bigr) \;=\; 1 + 3 + 8 \;=\; 12 \;=\; |V_{O_2}(3)|.
  \end{equation}

\item \textbf{Гауссово простое число Кулона $\widehat{M}(3) = 97$ и четырёхпетлевой физико-геометрический вывод постоянной тонкой структуры $\alpha_{\mathrm{fs}}^{-1}(d)$ в поле $\mathbb{Q}(\sqrt{2})$:}
  квадрат нормы комплексного импеданса элементарной ячейки $Z_{\mathrm{cell}}(d) = d_{\mathrm{in}}(\square) + i\,\mathrm{mod\_bits}(d) \in \mathbb{Z}[i]$ (объединяющего $d_{\mathrm{in}}(\square) = 4$ ребра квадратного плакета и $\mathrm{mod\_bits}(d) = d^2 = 9$ внутренних осевых переходов) равен
  \begin{equation}
    \label{eq:mhat-97-def}
    \widehat{M}(d) \;=\; d_{\mathrm{in}}(\square)^2 + \mathrm{mod\_bits}(d)^2 \;=\; 4^2 + 9^2 \;=\; 97,
  \end{equation}
  является простым числом в $\mathbb{Z}$ и в $\mathbb{Z}[i]$ исключительно при $d = 3$, а в сумме с полушириной подпороговой зоны Гейзенберга $\Delta\varphi_{\mathrm{disc}}(3) - 1 = d_{\mathrm{out}}(E)\,N_{\mathrm{tri}}(3) = 40$ даёт целую часть обратной постоянной тонкой структуры:
  \begin{equation}
    \label{eq:alpha-floor-137}
    \bigl\lfloor\alpha_{\mathrm{fs}}^{-1}(3)\bigr\rfloor \;=\; \widehat{M}(3) + \bigl(\Delta\varphi_{\mathrm{disc}}(3) - 1\bigr) \;=\; 97 + 40 \;=\; 137 \;=\; Q_\square(3) + \mathrm{mod\_bits}(3) \;=\; 11^2 + 4^2.
  \end{equation}
  На оболочке кубооктаэдра $O_2(d)$ с длиной ребра $|e|_2 = \sqrt{d-1} = \sqrt{2}$ и $N_{\mathrm{tri}}(d) - 1 = 7$ открытыми октантами вокруг источника ($1$ октант занят источником) обратная константа связи $\alpha_{\mathrm{fs}}^{-1}(d) = T_{\mathrm{vac}}(d) / \Gamma_{\mathrm{em}}(d)$ определяется отношением коэффициента пропускания вакуума $T_{\mathrm{vac}}(d)$ к двухшаговой амплитуде вершинного сцепления $\Gamma_{\mathrm{em}}(d)$ по четырём последовательным степеням $\widehat{M}(d)^{-1}, \dots, \widehat{M}(d)^{-4}$:
  \begin{align}
    \Gamma_{\mathrm{em}}(d) &= \frac{\sqrt{2}}{2\,\widehat{M}(d)} \;+\; \frac{1}{2\bigl(N_{\mathrm{tri}}(d) - 1\bigr)\widehat{M}(d)^2}, \label{eq:alpha-gamma-denom} \\
    T_{\mathrm{vac}}(d) &= 1 \;-\; \frac{\sqrt{2}}{d_{\mathrm{in}}(\square)\bigl(N_{\mathrm{tri}}(d) - 1\bigr)\widehat{M}(d)^3} \;-\; \frac{N_{\mathrm{tri}}(d)}{\bigl(N_{\mathrm{tri}}(d) - 1\bigr)^2\widehat{M}(d)^4}. \label{eq:alpha-tvac-numer}
  \end{align}
  Здесь в~\eqref{eq:alpha-gamma-denom} член порядка $\widehat{M}^{-1}$ есть одношаговая проекция единичного кванта вдоль ребра длины $\sqrt{2}$ с направляющим косинусом $1/\sqrt{2} = \sqrt{2}/2$, а член порядка $\widehat{M}^{-2}$ --- двухшаговый симметричный возвратный поток через $N_{\mathrm{tri}}(d) - 1 = 7$ открытых октантов; в~\eqref{eq:alpha-tvac-numer} из единичной проницаемости вакуума вычитаются трёхшаговая замкнутая петля по треугольной грани (распределённая по $d_{\mathrm{in}}(\square) = 4$ сторонам плакета и $7$ открытым октантам, порядок $\widehat{M}^{-3}$) и четырёхшаговая замкнутая петля по квадратной грани (сцепляющая все $N_{\mathrm{tri}}(d) = 8$ октантов через $(N_{\mathrm{tri}}(d) - 1)^2 = 49$ упорядоченных пар открытых октантов, порядок $\widehat{M}^{-4}$).

  Умножение числителя~\eqref{eq:alpha-tvac-numer} и знаменателя~\eqref{eq:alpha-gamma-denom} на общий целый знаменатель $d_{\mathrm{in}}(\square)(N_{\mathrm{tri}}(d) - 1)^2\widehat{M}(d)^4$ даёт точное представление $\alpha_{\mathrm{fs}}^{-1}(d)$ в поле $\mathbb{Q}(\sqrt{2})$ без трансцендентного числа~$\pi$:
  \begin{equation}
    \label{eq:alpha-fs-137}
    \alpha_{\mathrm{fs}}^{-1}(d) \;=\; \frac{T_{\mathrm{vac}}(d)}{\Gamma_{\mathrm{em}}(d)} \;=\; \frac{A_{\mathrm{den}}(d) + B_{\mathrm{den}}(d)\sqrt{2}}{A_{\mathrm{num}}(d) + B_{\mathrm{num}}(d)\sqrt{2}} \;=\; \frac{P_{\mathrm{rat}}(d) + Q_{\mathrm{rat}}(d)\sqrt{2}}{D_{\mathrm{rat}}(d)},
  \end{equation}
  где целочисленные коэффициенты при $d = 3$ равны:
  \begin{align}
    A_{\mathrm{num}}(3) &= \frac{d_{\mathrm{in}}(\square)}{2}\bigl(N_{\mathrm{tri}}(3) - 1\bigr)\widehat{M}(3)^2 \;=\; 131\,726, \label{eq:alpha-anum} \\
    B_{\mathrm{num}}(3) &= \frac{d_{\mathrm{in}}(\square)}{2}\bigl(N_{\mathrm{tri}}(3) - 1\bigr)^2\widehat{M}(3)^3 \;=\; 89\,441\,954, \label{eq:alpha-bnum} \\
    A_{\mathrm{den}}(3) &= d_{\mathrm{in}}(\square)\Bigl(\bigl(N_{\mathrm{tri}}(3) - 1\bigr)^2\widehat{M}(3)^4 - N_{\mathrm{tri}}(3)\Bigr) \;=\; 17\,351\,739\,044, \label{eq:alpha-aden} \\
    B_{\mathrm{den}}(3) &= -\bigl(N_{\mathrm{tri}}(3) - 1\bigr)\widehat{M}(3) \;=\; -679, \label{eq:alpha-bden}
  \end{align}
  что даёт десятичное значение
  \begin{equation}
    \label{eq:alpha-numeric-codata}
    \alpha_{\mathrm{fs}}^{-1}(3) \;=\; \frac{17\,351\,739\,044 - 679\sqrt{2}}{131\,726 + 89\,441\,954\sqrt{2}} \;=\; 137{,}035\,999\,187\dots
  \end{equation}
  (экспериментальное значение CODATA-2022: $137{,}035\,999\,177(21)$, расхождение $0{,}47$ стандартного отклонения).

\item \textbf{Жордановы замки $\mathrm{SL}_2(\mathbb{Z})$ и подъём на кривую рода $g(6) = 10$:}
  целочисленность нормированных матриц перехода $\widetilde{P}_k = I_2 \pm k N_0 \in \mathrm{SL}_2(\mathbb{Z})$ ($N_0^2 = 0$, \texttt{thm6}, \texttt{thm15}) выделяет четыре жордановых замка $k \in \{6, 12, 15, 40\}$, фиксируя границу Гейзенберга $40 \leftrightarrow 41$ и поднимая сцепленные поверхностно-объёмные состояния ($k = 2 \leftrightarrow k = 3$) на кривую $6$-й степени рода $g(6) = 10 \ge 2$ ($\chi(\mathcal{C}_6) = -18 < 0$).
\end{enumerate}
\end{theorem}

\begin{proof}
В пункте~(1) перебор подмножеств $13$-вершинного графа $B_1(3)$ даёт $13$ вершин ($K_1$), $12$ радиальных и $24$ поверхностных ребра ($K_2 = 36$), $24$ радиальных треугольника и $8$ треугольных граней оболочки ($K_3 = 32$) и $8$ тетраэдрических $4$-клик ($K_4 = 8$, по одной в каждом октанте), откуда $\chi(\mathcal{C}_{13}) = 13 - 36 + 32 - 8 = 1$. Прямое перемножение $3\times 3$ матриц~\eqref{eq:su2-generators-matrices} даёт
\begin{equation*}
  L_1 L_2 - L_2 L_1 \;=\; \begin{pmatrix} 0 & 0 & 0 \\ 1 & 0 & 0 \\ 0 & 0 & 0 \end{pmatrix} - \begin{pmatrix} 0 & 1 & 0 \\ 0 & 0 & 0 \\ 0 & 0 & 0 \end{pmatrix} \;=\; \begin{pmatrix} 0 & -1 & 0 \\ 1 & 0 & 0 \\ 0 & 0 & 0 \end{pmatrix} \;=\; L_3,
\end{equation*}
и циклически $[L_2, L_3] = L_1$, $[L_3, L_1] = L_2$. Для матричных единиц $(E_{ab})_{ij} = \delta_{ai}\delta_{bj}$ имеем $E_{ab} E_{cd} = \delta_{bc} E_{ad}$, откуда $[E_{ab}, E_{bc}] = E_{ac}$ при $a \ne c$ и $[E_{ab}, E_{ba}] = E_{aa} - E_{bb}$, что доказывает~\eqref{eq:su3-comm-step}--\eqref{eq:su3-comm-cartan}.

В пункте~(2) умножение каждого слагаемого в~\eqref{eq:alpha-gamma-denom} и~\eqref{eq:alpha-tvac-numer} на $d_{\mathrm{in}}(\square)(N_{\mathrm{tri}}(d)-1)^2\widehat{M}(d)^4$ переводит $\Gamma_{\mathrm{em}}(d)$ в $A_{\mathrm{num}}(d) + B_{\mathrm{num}}(d)\sqrt{2}$, а $T_{\mathrm{vac}}(d)$ --- в $A_{\mathrm{den}}(d) + B_{\mathrm{den}}(d)\sqrt{2}$ с коэффициентами~\eqref{eq:alpha-anum}--\eqref{eq:alpha-bden}. Умножение числителя и знаменателя~\eqref{eq:alpha-fs-137} на сопряжённый элемент $A_{\mathrm{num}} - B_{\mathrm{num}}\sqrt{2}$ в поле $\mathbb{Q}(\sqrt{2})$ даёт рационализованную форму $(P_{\mathrm{rat}} + Q_{\mathrm{rat}}\sqrt{2})/D_{\mathrm{rat}}$ со знаменателем Пелля $D_{\mathrm{rat}} = 2 B_{\mathrm{num}}^2 - A_{\mathrm{num}}^2 > 0$.
\end{proof}

\section{Двусортное проективно-ретракционное соответствие носителей \texorpdfstring{$\langle\mathcal{M},\mathcal{T}\rangle$}{<M,T>}: Трёхзаконная целочисленная механика Ньютона на \texorpdfstring{$O_2(3)\times\mathbb{Z}_{512}$}{O2(3) x Z512} и пять следствий дискретности решётки}
\label{sec:newtonian-homeomorphism}

Центральный структурный вывод Теорем~\ref{thm:d3-uniqueness}--\ref{thm:sm-alpha} состоит в том, что \textbf{все разделы фундаментальной физики, химии, молекулярной биологии и инженерии сводятся через двусортное проективно-ретракционное соответствие носителей $\langle\mathcal{M},\mathcal{T};\,\pi_{M\to T},\,\iota_{T\to M}\rangle$ к классической механике трёх законов Ньютона в точных целых числах на решётке кубооктаэдра $O_2(3)\times\mathbb{Z}_{512}$}:
\begin{enumerate}[label=\textbf{Закон Ньютона \Roman*.},leftmargin=*]
\item \textbf{Закон инерции в подпороговом зазоре зацепления ($|\Delta\varphi|\le 40 < 41 \implies \Phi_{\mathrm{rad}} = 0$):}
  при нулевом надпороговом моменте $\Delta^2\varphi(t) = 0 \implies \Delta\varphi(t) = v_0 = \mathrm{const}$ и $\varphi(t) = \varphi_0 + v_0 t$ (\texttt{thm1}, \texttt{thm36}). Стационарные орбиты в атоме, сверхпроводимость, баллистические экситоны и обратимые логические вентили Канала~1 движутся по инерции без излучения и без тепла Ландауэра ($\hat{n}_E = 0$).
\item \textbf{Симплектический закон динамики второго порядка с переменной инерцией узла ($\Delta^2\varphi = \lfloor N_{\mathrm{torque}}/D(\mathbf{x})\rfloor$):}
  единственная обратимая во времени рекурсия~\eqref{eq:verlet-law} есть второй закон Ньютона $a = \lfloor F/m\rfloor$ на $\mathbb{Z}_{512}$, где импеданс узла $D(\mathbf{x})\in[8\,192,\; 833\times 8\,192]$ играет роль локального момента инерции.
\item \textbf{Равенство действия и противодействия по всем 12 валам кубооктаэдра ($\mathcal{S}_3(-\Delta\varphi) = -\mathcal{S}_3(\Delta\varphi)$):}
  в силу нечётности синус-проекции (\texttt{odd\_holonomy\_current\_pairwise\_conservation\_all\_512}) парные потоки между любыми двумя соседними узлами в сумме тождественно равны нулю ($F_{\mathbf{x}\to\mathbf{y}} + F_{\mathbf{y}\to\mathbf{x}} \equiv 0$), напрямую обеспечивая сохранение импульса, момента импульса, заряда и тождества $d_1\circ d_0 \equiv 0$, $d_2\circ d_1 \equiv 0$.
\end{enumerate}

Все эффекты, для которых в континуальных теориях XX~века (НКМ, СТО, ОТО, КТП) вводились отдельные постулаты, возникают не из отказа от трёх законов Ньютона, а из \textbf{пяти механических свойств самой дискретной решётки $O_2(3)\times\mathbb{Z}_{512}$}, сведённых в Таблицу~\ref{tab:newtonian-homeomorphism}.

\begin{table}[htbp]
\centering
\caption{Двусортное проективно-ретракционное соответствие носителей $\langle\mathcal{M},\mathcal{T}\rangle$: сведение континуальных физических и инженерных постулатов к трёхзаконной целочисленной механике Ньютона на $O_2(3)\times\mathbb{Z}_{512}$.}
\label{tab:newtonian-homeomorphism}
\footnotesize
\begin{tabular}{@{}p{0.20\textwidth}p{0.36\textwidth}p{0.39\textwidth}@{}}
\toprule
\textbf{Раздел и свойство решётки} & \textbf{Континуальная идеализация на $\mathcal{T}$-слое ($\mathbb{R}^{3,1}$)} & \textbf{Дискретная ньютоновская механика на $\mathcal{M}$-слое ($O_2(3)\times\mathbb{Z}_{512}$)} \\
\midrule
\textbf{1. Кинематика и релятивизм}\newline ($1$ ребро $h_L$ за такт $h_T$, $c_0=\sqrt{2}c$) &
Постулирует сокращение пространства-времени Минковского и трактует запутанность как нелокальное дальнодействие. &
Жёсткая ньютоновская решётка $\mathbb{Z}^3$ и такты $t\in\mathbb{Z}$; потолок скорости по ребру $c_0 = h_L/h_T = \sqrt{2}c$ замедляет внутренние роторы движущихся частиц и передаёт фазу запутанности по рёбрам со скоростью $c_0$ (граница Цирельсона $2\sqrt{2}$). \\
\addlinespace[3pt]
\textbf{2. Гравитация и космология}\newline ($D(\mathbf{x})\in[\varepsilon,833\varepsilon]$, фактор $16/3$) &
Искривляет пустой континуум ($G_{\mu\nu}=8\pi T_{\mu\nu}$), порождая сингулярности чёрных дыр и требуя $85\%$ тёмной материи. &
Второй закон Ньютона $\Delta^2\varphi = \lfloor N/D(\mathbf{x})\rfloor$ в среде с переменной инерцией; потолок Паули $D_{\max}=833\varepsilon$ исключает сингулярности, а вороновское гало $16/3$ даёт плоские кривые вращения галактик без тёмной материи. \\
\addlinespace[3pt]
\textbf{3. Квантовая механика}\newline (Люфт зацепления зуба $\Delta\varphi_{\mathrm{disc}}=41$) &
Заменяет траектории волнами вероятности в гильбертовом пространстве, коллапсом при измерении и теплом $k_{\mathrm{B}}T\ln 2$. &
Детерминированная траектория Верле в каждой ячейке; целочисленный порог $\hat{n}_E = \lfloor|\Delta\varphi|/41\rfloor$ даёт нулевое излучение при $|\Delta\varphi|\le 40$ (инерция по Закону~I) и правило Борна $P(\mu)=\mu^2/512$ при $72^2=5\,184$ скрытых состояниях. \\
\addlinespace[3pt]
\textbf{4. Электродинамика и КТП}\newline (Кольцо $\mathbb{Z}_{512}$, упор Паули $256$) &
Постулирует калибровочные расслоения и требует УФ-перенормировки расходящихся петлевых интегралов. &
Проскок на оборот $\Delta\varphi = 512 q + r$ даёт квантованный вихревой поток $(d_1 r)_f = 512 n_f$ и законы Максвелла ($d^2\equiv 0$); противофаза $256$ создаёт упор Паули с жёстким УФ-потолком $\hat{n}_E\le 6$. \\
\addlinespace[3pt]
\textbf{5. Химия, биология и инженерия}\newline ($12$ валов $O_2(3)$, седловина $Q_\square=128$) &
Опирается на эмпирические правила орбиталей, уравнения Навье--Стокса, предел Шокли--Квейссера $33\%$ и коэффициенты усталости. &
Слои $S_k=10k^2+2$ и $\mathcal{C}_{13}$ порождают таблицу Менделеева; седловина $Q_\square=128$ ($\mathcal{C}_3=0$) даёт катализ $587\times$, $C_4$-шлюз $\mathrm{RyR1}$ и бесфононную кремниевую ячейку $120+8=128$; $12$-лучевой расщепитель ($256/12=21<41$) выпрямляет тепло без памяти. \\
\bottomrule
\end{tabular}
\end{table}

\begin{theorem}[Сокращение масштаба $A = K\cdot h_L$, целая размерность Хаусдорфа $D_H = 3$ вакуумного носителя Вороного и происхождение фракталов природы $D_H = \ln N_{\mathrm{sub}}/\ln K < 3$]
\label{thm:scale-cancellation-fractals}
При переходе к любому когерентному макроэтажу с целочисленным шагом ребра $A = K\cdot h_L$ ($K\in\mathbb{N}$, \texttt{scale\_invariance\_of\_holonomy\_and\_gate\_operators}, \texttt{scale\_cancellation\_voronoi\_hausdorff\_and\_fractal\_subcascades} и \texttt{cross\_floor\_shell\_closure\_isomorphism\_operator}):
\begin{enumerate}[label=\textbf{(\arabic*)},leftmargin=*]
\item \textbf{Тождественное сокращение $K^2$ во втором законе Ньютона и инвариантность кванта действия $s_{\mathrm{floor}}(K) = s_0$:}
  при перемасштабировании этажа $(\mu,\Sigma\mu)\mapsto(K\mu,K\Sigma\mu)$ и $(\zeta_r,\zeta_i)\mapsto(K^2\zeta_r,K^2\zeta_i)$ числитель момента и знаменатель инерции в $\Delta^2\varphi = \lfloor N_{\mathrm{torque}}/D(\mathbf{x})\rfloor$ однородны степени~$2$ по~$K$:
  \begin{equation}
    \label{eq:scale-k2-cancel}
    N_{\mathrm{torque}}^{(K)}\cdot D^{(1)} \;\equiv\; N_{\mathrm{torque}}^{(1)}\cdot D^{(K)},
  \end{equation}
  а сопряжённый шаг импульса $\Delta p(K) = p_0/K$ сохраняет квант действия $s_{\mathrm{floor}}(K) = \Delta p(K)\cdot(K h_L) = s_0 = \hbar/2$ и 9-битное кольцо $\mathbb{Z}_{512}$ на всех этажах масштаба~$K$.
\item \textbf{Сплошное самоподобие ячеек Вороного вакуумного носителя ($D_H = d = 3$):}
  макро-ячейка Вигнера--Зейтца (ромбический додекаэдр с ребром $A = K\cdot h_L$) содержит $N_{\mathrm{pack}}(d,K) = K^d$ элементарных ячеек и удовлетворяет мультипликативному самоподобию $N_{\mathrm{pack}}(d,K_1 K_2) = N_{\mathrm{pack}}(d,K_1)\,N_{\mathrm{pack}}(d,K_2)$ с целой размерностью Хаусдорфа--Морана $D_H^{(\mathrm{bulk})} = \ln(K^3)/\ln K = 3$, сохраняя теорему Гаусса и закон $1/R^2$ в вакууме.
\item \textbf{Происхождение природных фракталов как подкубических каскадов отбора ($1 < N_{\mathrm{sub}} < K^3 \implies 0 < D_H = \ln N_{\mathrm{sub}}/\ln K < 3$):}
  когда локальные правила отбора автомата оставляют на шаге масштаба~$K$ активное подмножество $N_{\mathrm{sub}} < K^3$ (с непустым вырезанным дополнением $K^3 - N_{\mathrm{sub}} > 0$), итерация по $m$ вложенным этажам порождает пористый дискретный фрактал с $N_{\mathrm{sub}}^m < K^{3m}$ узлами и дробной размерностью Хаусдорфа--Морана $D_H = \ln N_{\mathrm{sub}}/\ln K < 3$:
  \begin{itemize}[leftmargin=*]
  \item \emph{Октантная губка Серпинского--Поста в Канале~2 и 3-SAT ($K = 2$):} $N_{\mathrm{sub}} = 2^3 - 1 = 7 < 8 = 2^3$, что даёт $D_H = \ln 7/\ln 2 \approx 2{,}807355$;
  \item \emph{Кубооктаэдрическая губка $\mathcal{C}_{13}$ в кубе Мура $3\times 3\times 3$ ($K = d = 3$):} каузальность одного такта сохраняет $N_{\mathrm{sub}} = N_\varphi(3) = 13 < 27 = 3^3$ узлов, давая $D_H = \ln 13/\ln 3 \approx 2{,}334718$;
  \item \emph{Иерархическая кластеризация связанной материи на лестнице тонкой структуры ($K = \lfloor\alpha_{\mathrm{fs}}^{-1}\rfloor = 137$):} замкнутые оболочечные кластеры $N_{\mathrm{sub}}\in\{13, 42, 147\} \ll 137^3$ дают $D_H = \ln N_{\mathrm{sub}}/\ln 137 \in [0{,}522;\,0{,}999] < 3$;
  \item \emph{Фрактальные сосудистые и бронхиальные деревья живых организмов (\texttt{allometryExponentRatioOp}):} $4$-кратное ветвление по квадратной грани кубооктаэдра $k = d_{\mathrm{in}}(\square) = 4$ в $3$-мерном пространстве даёт показатель метаболической аллометрии Клейбера $(d_{\mathrm{in}}(\square)-\mu_{\mathrm{vac}})/d_{\mathrm{in}}(\square) = 3/4$.
  \end{itemize}
\item \textbf{Операторы логарифма $\log_a(b)$ и рациональной степени/корня $a^{p/q}$ на $\mathbb{Z}_{4Q}$ и обращение размерностей Хаусдорфа--Морана $K^{D_H} = N_{\mathrm{sub}}$ (\texttt{rotatorLog2HomogQOp}, \texttt{rotatorLogBaseRatQOp}, \texttt{rotatorPowRatQOp}):}
  для любого четверть-периода $Q\ge 1$ на кольце $\mathbb{Z}_{4Q}$ разложение $x = 2^k + m$ на характеристику $k = \lfloor\log_2 x\rfloor$ и фазу мантиссы первого квадранта $\mu = \lfloor m Q/2^k\rfloor\in[0,Q]$ ($\mathcal{C}_Q(\mu)=Q-\mu$, $\mathcal{S}_Q(\mu)=\mu$) задаёт однородный числитель двоичного логарифма
  \begin{equation}
    \label{eq:log2-homog-q}
    \operatorname{Log2}_{d,Q}(k,\mu) \;:=\; \bigl(d_{\mathrm{in}}(\square)\,N_{\mathrm{tri}}(d)\bigr)Q^2\bigl(k Q + \mathcal{S}_Q(\mu)\bigr) + \mathcal{C}_Q(\mu)\mathcal{S}_Q(\mu)\Bigl(\bigl(|V_{O_2}(d)|-\mu_{\mathrm{vac}}(d)\bigr)Q + d\bigl(\mathcal{C}_Q(\mu)-\mathcal{S}_Q(\mu)\bigr)\Bigr)
  \end{equation}
  и обратный числитель мантиссы экспоненты
  \begin{equation}
    \label{eq:exp2-homog-q}
    \operatorname{Exp2Mant}_{d,Q}(r) \;:=\; \bigl(d_{\mathrm{in}}(\square)\,N_{\mathrm{tri}}(d)\bigr)Q^2\mathcal{S}_Q(r) - \mathcal{C}_Q(r)\mathcal{S}_Q(r)\Bigl(\bigl(|V_{O_2}(d)|-\mu_{\mathrm{vac}}(d)\bigr)Q - d\bigl(\mathcal{C}_Q(r)-\mathcal{S}_Q(r)\bigr)\Bigr),
  \end{equation}
  удовлетворяющие кубической однородности $\operatorname{Log2}_{d,KQ}(k,K\mu) \equiv K^3\,\operatorname{Log2}_{d,Q}(k,\mu)$ и $\operatorname{Exp2Mant}_{d,KQ}(K r) \equiv K^3\,\operatorname{Exp2Mant}_{d,Q}(r)$. Благодаря этому масштаб $Q^3$ тождественно сокращается в отношении $\operatorname{LogBaseRat}_{d,Q}(a,b) = \operatorname{reduceRat}\bigl(\operatorname{Log2Rot}_{d,Q}(b),\,\operatorname{Log2Rot}_{d,Q}(a)\bigr)$, а композиция $\operatorname{Exp2Scaled}_{d,Q}(y) = \bigl(Q + \operatorname{Exp2Mant}_{d,Q}(y\bmod Q)\bigr)\,2^{\lfloor y/Q\rfloor}$ при целочисленном аргументе $y = \operatorname{round}\bigl(p\,\operatorname{Log2Rot}_{d,Q}(a)/q\bigr)$ задаёт оператор произвольной рациональной степени и корня $\operatorname{PowRat}_{d,Q}(a,p,q) \approx a^{p/q}$. При $d=3$ и $Q=128$ эти операторы вычисляют $\log_2(8)=(3,1)$, $\log_3(9)=\log_{10}(100)=(2,1)$, $\log_2(3)=(203,128)$, $\log_2(7)=(359,128)$, $\log_3(13)=(474,203)$, $\log_{10}(2)=(128,425)$, точные степени и корни $2^3=(8,1)$, $3^3=9^{3/2}=(27,1)$, $\sqrt{4}=\sqrt[3]{8}=(2,1)$, $\sqrt{9}=\sqrt[3]{27}=(3,1)$, $\sqrt{16}=8^{2/3}=(4,1)$, аллометрию Клейбера $16^{3/4}=(8,1)$, подходящую дробь Пелля $\sqrt{2}=2^{1/2}=(181,128)$, $\sqrt{3}=3^{1/2}=(111,64)$ и обращают $K^{D_H} = N_{\mathrm{sub}}$: $2^{203/128}=(3,1)$, $\operatorname{round}(2^{359/128})=7$ и $3^{474/203}=(13,1)$.
\end{enumerate}
\end{theorem}

\begin{proof}
В пункте~(1) при замене $\mu \mapsto K\mu$ каждая квадратичная форма по амплитудам умножается на $K^2$, откуда $N_{\mathrm{torque}}^{(K)} = K^2 N_{\mathrm{torque}}^{(1)}$ и $D^{(K)} = K^2 D^{(1)}$, что даёт~\eqref{eq:scale-k2-cancel}. В пункте~(2) объём ромбического додекаэдра со стороной $K h_L$ равен $K^3 V_{\mathrm{WS}}(h_L)$, откуда $N_{\mathrm{pack}}(d, K_1 K_2) = (K_1 K_2)^d = K_1^d K_2^d$. В пункте~(4) при замене $(Q, \mu) \mapsto (KQ, K\mu)$ имеем $\mathcal{C}_{KQ}(K\mu) = K\mathcal{C}_Q(\mu)$ и $\mathcal{S}_{KQ}(K\mu) = K\mathcal{S}_Q(\mu)$, поэтому каждое слагаемое в~\eqref{eq:log2-homog-q} и~\eqref{eq:exp2-homog-q} является однородным многочленом третьей степени от $(Q, \mathcal{C}_Q, \mathcal{S}_Q)$ и приобретает общий множитель $K^3$, который сокращается в отношении $\operatorname{Log2Rot}_{d,Q}(b)/\operatorname{Log2Rot}_{d,Q}(a)$.
\end{proof}

\section{Формальная верификация в Lean~4}
Все теоремы настоящей работы формально верифицированы в пакете \texttt{lean/MtCarrier/} в репозиториях \url{https://github.com/AI-Guiders/discrete-quantum-mechanics} и \url{https://github.com/AI-Guiders/scientific-papers} (\texttt{0 sorryAx}, \texttt{0 Classical.choice}).

\begin{thebibliography}{9}
\bibitem{Malcev1970}
А.\,И.~Мальцев, \emph{Алгебраические системы}, М.: Наука, 1970 (англ.\ пер.: Springer-Verlag, 1973).
\bibitem{LandauLifshitz}
Л.\,Д.~Ландау, Е.\,М.~Лифшиц, \emph{Теоретическая физика}, Т.~I--IV, М.: Наука / Физматлит.
\bibitem{Moura2021}
L.~de~Moura and S.~Ullrich, ``The Lean~4 theorem prover and programming language,'' in \emph{CADE~28}, LNCS vol.~12699, Springer, 2021, pp.~625--635.
\end{thebibliography}

\end{document}
